% !TEX root = ../paper.tex
\section{Introduction} \label{sec:introduction}

Local search is a fundamental paradigm in solving the satisfiability problem: find an assignment close in Hamming distance to the initial assignment that satisfies the formula, if one exists. Papadimitriou \cite{Papadimitriou91} was the first to employ this idea in a randomized poly-time 2-SAT algorithm. Sch{\" o}ning \cite{DSchoning02} showed that a slight modification of this algorithm yields a running time of $(2 - 2/k)^n$ for $k$-SAT. Dantsin et al. \cite{DantsinGHKKPRS02} considered a deterministic version of local search and gave a deterministic $(2 - 2/(k+1))^n$ time algorithm. Brueggemann and Kern \cite{BrueggemannK04} and Kutzkov and Scheder \cite{KS} improved this deterministic local search procedure and obtained faster deterministic 3-SAT algorithms. Moser and Scheder \cite{MoserS11} eventually considered a variant of the local search problem and used it to give a deterministic $k$-SAT algorithm matching the running time of Sch{\" o}ning's.

Despite the success of local search, the fastest known $k$-SAT algorithm PPSZ and its improvements follow a different approach: pick a random variable $x$, if its value is not easily seen to be forced\footnote{For example if $x$ appears in a unit clause, or if such a clause can be derived in small width resolution.} then assign it randomly and continue (see \cite{PaturiPZ99,PaturiPSZ05,Hertli14,DBLP:conf/stoc/HansenKZZ19,Scheder21}). The analysis of this simple yet powerful algorithm consists of a combinatorial theorem relating the size and the structure of the set of satisfying assignments of a $k$-CNF. The strength of this combinatorial theorem is further manifested by the fact that the state-of-the-art depth-3 circuit lower bounds are built on it \cite{PaturiPZ99,PaturiPSZ05}.

This curious interaction between lower bounds and algorithms has become less of a surprise over the years. Williams \cite{Williams13} initiated a whole new line of inquiry by showing that improved satisfiability algorithms for a circuit class automatically imply lower bounds for the same class. Conversely, almost all known circuit lower bound techniques have been adopted in satisfiability algorithms (see e.g. \cite{ImpagliazzoMP12,ChenKKSZ15}). Within this context the role of local search is unclear.

\begin{question}
Can we derive lower bounds from local search algorithms?
\end{question}

This question is also motivated by a lack of progress
in improving depth-3 circuit lower bounds and 
related upper bounds on $k$-SAT algorithms.

\bigskip
\noindent{\bf{Depth-3 circuit lower bounds.}} A $\Sigma^k_3$ circuit is an 
Or-And-Or circuit where the bottom fan-in is bounded by $k$, i.e., a disjunction of $k$-CNFs. We use $\Sigma^k_3(f)$ to denote the minimum number of $k$-CNFs in a $\Sigma^k_3$ circuit computing a function $f$. The study of 
these circuits was advocated by Valiant \cite{Valiant77} who showed that a 
strong enough $\Sigma^k_3$ lower bound for every fixed $k$ implies 
a super-linear lower bound for series-parallel circuits. Moreover, \cite{GKW21circuit} showed that strong lower bounds even for small constant $k$ such as $k=16$ imply various circuit lower bounds including better general circuit lower bounds.
The technique of \cite{PaturiPZ99} gives a lower bound of $\Omega(2^{n/k})$ 
for the parity function and it is known to be 
tight. In fact this results in a $\Omega(n^{\frac{1}{4}} 2^{\sqrt{n}})$ lower bound
for computing parity by unrestricted 
depth-3 circuits which is tight up to a constant factor. 
A further improvement comes from  \cite{PaturiPSZ05} which 
gives a lower bound of $2^{cn/k}$ where $c > 1$ for the BCH code. 
At this point, this is the best known
lower bound for computing any explicit function by $\Sigma^k_3$ circuits.
%The exact depth-3 complexity of the BCH code remains unknown.

\bigskip
%\noindent{\bf{Depth-3 complexity of Majority.}} 
Majority is a natural candidate 
for going beyond the $2^{\Omega(\sqrt{n})}$ depth-3 circuit lower bound.
The natural $\Sigma^k_3$ circuit for 
computing Majority has size $2^{O(n\log(k)/k)}$ (which implies an unrestricted depth-3 size upper bound of $2^{O(\sqrt{n\log n})}$). H\r{a}stad, Jukna and Pudl\'ak \cite{HastadJP95} introduced the intriguing notion of \emph{$k$-limits} to capture the depth-3 complexity of various functions  
and proved a lower bound of $2^{\Omega(\sqrt{n})}$ where the constant factor in the 
exponent is improved  over the constant one could obtain from Switching Lemma. 
Regarding $\Sigma^k_3$ circuits computing Majority, their result implies size lower bounds of  $1.414^n$ for $k = 2$, and $1.154^n$ for $k = 3$, and for $k \ge 4$ it yields nothing. The $\Sigma^2_3$ bound is known to be essentially tight \cite{PaturiSZ00}. More recently,
\cite{LecomteRT22} proved a tight lower bound of $2^{\Omega(n\log(k)/k)}$ for 
computing Majority by $\Sigma^k_3$ circuits where each 
And gate depends on at most $k$ variables.
Further, \cite{Amano23Majority} studied the effect of negations for $\Sigma^k_3$ circuits computing majority.
However, the question of proving tight lower bounds for computing Majority by depth-3 circuits 
(even for fan-in 3 circuits) remains open.

%The apparent barriers faced by PPSZ both in terms of 
%satisfiability as well as lower bounds, further justify this question.
%In fact these encoding arguments are so powerful that they constitute 
%the main component of the state-of-the-art depth-3 circuit lower bounds as well.


\bigskip
\noindent
{\bf $k$-SAT upper bounds:}
Lack of progress in improving the savings beyond $\Omega(\frac{1}{k})$
for $k$-SAT algorithms led researchers to consider
SSETH (Super Strong Exponential Time Hypothesis).
SSETH is the hypothesis that $k$-SAT cannot be solved with \emph{savings} asymptotically 
more than $1/k$, i.e., 
there is no $2^{(1 - \epsilon_k)n}$ time $k$-SAT 
algorithm with $\epsilon_k = \omega(1/k)$. 
%PPSZ and many other algorithms have savings 
%$\Omega(1/k)$. 
However, SSETH is known to be 
false on average \cite{DBLP:conf/sat/VyasW19,LincolnY20}, 
that is, 
the satisfiability of almost all $k$-CNFs can be decided with 
much larger savings. 
It is thus not unreasonable to attempt to get such savings 
even in the worst-case. 
Yet we cannot hope to achieve such a savings for 
a large subclass of PPSZ-style algorithms \cite{DBLP:conf/coco/SchederT20}.

\bigskip
It appears that making progress towards larger $k$-SAT savings as well as depth-3 circuit lower bounds requires new ideas. We argue that local search has the potential to achieve this goal, and as an evidence for this claim, we apply local search ideas to give a new $\Sigma^3_3$ lower bound for Majority function.


\subsection{Local enumeration, $k$-SAT and $\Sigma^k_3$ lower bounds for Majority function}
The local search problem is formally defined as follows.

\bigskip
\noindent{\bf{\ballsat{k}{t}}.} Given an $n$-variable $k$-CNF $F$, a parameter $t$, and an assignment $\alpha$, decide if there is a satisfying assignment $\beta$ of $F$ such that $d(\alpha, \beta) \le t(n)$, where $d(\cdot)$ is the Hamming distance.

\bigskip
Dantsin et al. \cite{DantsinGHKKPRS02} gave a simple branching algorithm which solves \ballsat{k}{t} in time $\poly(n) \cdot k^t$. This already gives a non-trivial algorithm for 3-SAT: solve \ballsat{k}{\frac{n}{2}} 
starting with the all-0  and all-1 assignments. To get a non-trivial algorithm for larger $k$, they used \emph{covering codes}, i.e., a small and asymptotically optimal number $C(n, t)$ of Hamming balls of a given radius $t$ that cover the entire $n$-dimensional Boolean cube. Then an upper bound of $C(n, r)\cdot k^t$ follows immediately for $k$-SAT by solving \ballsat{k}{t} starting with 
the centers of each of the balls in the covering code. Setting $t = \frac{n}{k + 1}$ minimizes this quantity. 
Thus improved upper bounds 
for \ballsat{k}{t} immediately imply 
improved $k$-SAT upper bounds, 
and indeed this is what \cite{BrueggemannK04,KS} did  by
proving an upper bound of $c^t$ for some $c < 3$ for $\ballsat{3}{t}$.
However, this improvement in local search is not
sufficient to yield better upper  
bounds for $k$-SAT when we want to use the technique for large $t$. It is conceivable that improved bounds for large $t$ combined with covering codes would yield improved $k$-SAT algorithms. This leads to the following question:

\begin{question}
What is the complexity of \ballsat{k}{\epsilon n} where $0 < \epsilon \le \frac{1}{2}$?
\end{question}

It is also natural to consider the enumeration problem for \ballsat{k}{t}: 
enumerate all satisfying assignments within Hamming distance $t$ of an 
initial assignment. We note that even a weaker form of this problem 
already captures the circuit complexity of Majority function. For 
this purpose, we introduce the following class of parameterized problems.

\bigskip
\noindent{\bf{\enumballsat{k}{t}}.} Given an $n$-variable $k$-CNF $F$ 
and an initial assignment $\alpha$, 
output all satisfying assignments of $F$ at a Hamming distance $t$
from $\alpha$ assuming that there are no satisfying assignments
of $F$ at a Hamming distance of less than $t$ from $\alpha$.

\bigskip
We observe that upper bounds on \enumballsat{k}{t} imply depth-3 circuit lower bounds and $k$-SAT algorithms.

\begin{proposition}
Assume that \enumballsat{k}{t} can be solved in randomized expected time $b(n, k, t)$. Then 
\begin{enumerate}
	\item any $\Sigma^k_3$ circuit requires at least $\binom{n}{n/2}/b(n, k, \frac{n}{2})$ size for computing Majority function.
	\item $k$-SAT can be solved in time $O(\sum_{t = 1}^{n/2}b(n, k, t))$.
\end{enumerate}
\end{proposition}


\begin{proof}
1) Consider a $\Sigma^k_3$ circuit $C$ that computes Majority function.
We will write $C= \bigvee_{i = 1}^m F_i$, where each $F_i$ is a $k$-CNF. 
None of the $F_i$s has a satisfying assignment with Hamming 
weight less than $n/2$. 
By assumption we can enumerate all satisfying assignments of Hamming weight exactly $n/2$ in expected time $b(n, k, \frac{n}{2})$. 
This in particular implies that the total number of such 
satisfying assignments for each $F_i$ is at most $b(n, k, \frac{n}{2})$. 
Since $F_i$s should together cover all 
assignments of Hamming weight exactly 
$n/2$ and since there are $\binom{n}{n/2}$ such assignments, 
the claim follows.

2) We can trivially check if there is a satisfying assignment of Hamming weight at most $n/2$  in $1 + \sum_{t = 1}^{n/2}b(n, k, t)$ steps. In the same number of steps we can check if there is a satisfying assignment of Hamming weight at least $n/2$.
\end{proof}

Observe that $b(n, k, \frac{n}{2})$ cannot be too small: 
define the $k$-CNF $\Maj_{n, k}$ by partitioning the $n$ variables into 
sets of size $2(k - 1)$ and 
by including all positive clauses of size $k$
from each of the parts. It is easy to see that every satisfying assignment 
of this formula must set at least $k - 1$ variables in each 
part to 1 and the total number satisfying assignments with 
Hamming weight $n/2$ is $2^{(1 - O(\log(k)/k))n}$, 
thus $b(n, k, \frac{n}{2}) \ge 2^{(1 - O(\log(k)/k))n}$. It 
follows that a matching upper bound 
for \enumballsat{k}{\frac{n}{2}} refutes SSETH and 
gives $\Sigma^k_3$ lower bound of $2^{\Omega(\frac{\log k}{k}n)}$ 
for Majority which in turn implies a $2^{\omega(\sqrt{n})}$ unrestricted depth-3 circuit lower bound, breaking a decades long barrier.


\subsection{Our contributions}

In this paper, we study \enumballsat{3}{\epsilon n} 
and obtain new algorithms and lower bounds. Note that this is the first non-trivial instance of \enumballsat{k}{t}, since \enumballsat{2}{t} can be solved in $2^t$ steps using a simple extension of the local search algorithm, and it is easy to see that this is tight for $t \le n/2$ by considering the 2-CNF consisting of $t$ disjoint monotone clauses.

\begin{theorem}[Main result]
\label{thm:main-result}
\enumballsat{3}{t} can be solved in expected time
%$\left(\frac{3}{2+1/\sqrt{3}}\right)^n\left(\frac{2 + 1/\sqrt{3}}{9}\right)^t

\begin{enumerate}
	\item $3^t$, for $t \le \frac{n}{3}$,
	\item $1.164^n \times 1.9023^t$, for $\frac{n}{3} < t \le \frac{3n}{7}$,
	\item $1.1962^n \times 1.7851^t$, for $\frac{3n}{7} < t \le \frac{n}{2}$.
\end{enumerate}

In particular, \enumballsat{3}{\frac{n}{2}} can be solved in expected time $1.598^n$. 
\end{theorem}

Consequently, we get
\begin{corollary}
$\Sigma^3_3(\Maj) \ge 1.251^{n - o(n)}$.
\end{corollary}

Our lower bound is the best known bound
to compute Majority function by $\Sigma^3_3$ circuits.
Note that $\Maj_{n, 3}$ has $6^{n/4} \simeq 1.565^n$ satisfying assignments of Hamming weight $n/2$. Our bound is not too far from the optimal bound 
and it is a substantial improvement over $3^{n/2} \simeq 1.732^n$.

In the following, we explain how our approach to enumeration differs 
from the well-known approaches.
The $k^t$ algorithm used by \cite{DantsinGHKKPRS02} which 
solves \ballsat{k}{t} is a simple branching procedure. Without 
loss of generality assume that the initial assignment is all-0. If 
there is no monotone clause in the formula, 
then the all-0 assignment satisfies the formula. 
Otherwise select a monotone clause $C = x_1 \vee \ldots \vee x_k$. 
Then for each $x_i$, recursively solve \ballsat{k}{(t-1)} 
for the formula restricted by $x_i = 1$. 
An obvious weakness of this algorithm is that 
if the depth of the recursion tree is more than $n/k$, 
then some assignments will be considered more than once in 
the tree thus leading to redundant computation.

\bigskip
%\noindent{\bf Randomized Monien-Speckenmeyer.} 
Our starting point is to search
the recursion tree 
(which we call a transversal tree following the hypergraph nomenclature) 
so each satisfying assignment (which we call a transversal in the following) within the ball is visited exactly once.
We observe that such a non-redundant search can be conducted 
with any clause ordering and any fixed ordering of variables 
within the clause. During the search of the transversal tree, 
it is easy to decide whether 
a subtree contains any new satisfying assignments by considering 
the labels of the child edges of the nodes along the path.
If there are no new satisfying assignments in a subtree,  we will prune 
it.
It turns out that this approach is isomorphic to the 
seminal method of 
Monien-Speckenmeyer \cite{MonienS85} where for each $i$ we 
recursively solve the problem under the 
restriction $x_1 = \ldots = x_{i - 1} = 0, x_i = 1$. 
However, it is not clear
how to improve 
upon the bound obtained by \cite{MonienS85}.
We show that by choosing 
the clause ordering  carefully and randomly ordering 
clause variables, a better bound can be obtained. 
In other words, we will consider 
\emph{randomized} Monien-Speckenmeyer trees with careful clause ordering.
The crux of our contribution is a  new 
analysis of randomized transversal trees.


\bigskip
\noindent{\bf{Connection to hypergraph Tur{\'a}n problems.}} 
Recall that a transversal in a hypergraph is a set of 
vertices that intersects every hyperedge. 
The recent work \cite{FranklGT22} gives a connection 
between depth-3 circuits and 
transversals\footnote{\cite{FranklGT22} results 
are stated in terms of cliques which are dual to transversals.}. 
Here we find 
another connection. Given $n$, $t$, and $k$, 
let $R^+(n, t, k)$ be the maximum number of 
transversals of size $t$ in an $n$-vertex $k$-graph 
with no transversal of size $t - 1$. 
Since a $k$-graph can be viewed as a monotone $k$-CNF, our algorithm 
enumerates minimum size transversals and thus gives new upper 
bounds for $R^+(n, t, k)$. 
Mantel's theorem, which is a special case of 
Tur{\'a}n's seminal theorem, gives the 
exact value $R^+(n, n - 2, 2) = n^2/4$, 
and the Tur{\'a}n problem for 
3-graphs can be phrased as 
showing $R^+(n, n - 3, 3) = (5/9 + o(1))n^3$ 
(see \cite{MR2866732} for a thorough survey of this 
and related problems). Our technique allows us to derive 
new bounds for $R^+(n, t, 3)$, 
where $t = \Theta(n)$. 
Although we currently do not get any useful results 
for  $t = n - o(n)$, we hope that our techniques can be extended to make progress for
this regime of parameters.


\bigskip
\noindent{\bf{Enumeration algorithms for CNFs with bounded negations.}} 
As an additional application of our enumeration techniques, we get the following enumeration algorithm:
\begin{theorem}
\label{thm:enum3negate}
Let $F$ be a CNF of arbitrary width where either each clause contains at most $3$ negative literals or each clause contains at most $3$ positive literals. Then, we can enumerate all minimal satisfiable solutions of $F$ in time $O(1.8204^n)$.
\end{theorem}
% As far as we know, no non trivial satisfiability algorithm let alone an enumeration algorithm was previously known for such CNFs.






















